% Tipo de documento. En este caso es un art�culo, para folios A4, tama�o de la fuente 11pt y con p�gina separada para el t�tulo
\documentclass[a4paper,11pt,titlepage]{article}

% Carga de paquetes necesarios. OrdenesArticle es un paquete personalizado
\usepackage[spanish]{babel} 
%quitado temporalmente -> \RequirePackage[T1]{fontenc}
\RequirePackage[ansinew]{inputenx} 
\usepackage[spanish,cap,cont,title,fancy]{OrdenesArticle}
\usepackage{array}
\usepackage{graphicx}
\usepackage{hyperref}
\usepackage{pifont}
\usepackage{listings}
\usepackage[usenames,dvipsnames]{color}
\usepackage{colortbl}
\usepackage{color}
\usepackage{ifthen}
\usepackage{longtable}
\hypersetup{bookmarksopen,bookmarksopenlevel=4,linktocpage,colorlinks,urlcolor=blue,citecolor=blue,
						linkcolor=blue,filecolor=blue,pdfnewwindow,
						pdftitle={Funcionalidad 3},
						pdfauthor={Juan Andrada Romero, Juan Jos� Antequera Flores, Jose Domingo L�pez L�pez},
						pdfsubject={Modelos Avanzados de Bases de Datos}}


% Macro para definir una lista personalizada 
\newenvironment{milista}%
{\begin{list}{\textbullet}%
{\settowidth{\labelwidth}{\textbullet} \setlength{\leftmargin}{\dimexpr\labelsep+\labelwidth+5pt}
\setlength{\itemsep}{\dimexpr 0.5ex plus 0.25ex minus 0.25ex}
\setlength{\parsep}{\itemsep}
\setlength{\partopsep}{\itemsep}
\addtolength{\topsep}{-7.5pt}
}}%
{\end{list}}

% Macro para insertar una imagen
%       Uso: \imagen{nombreFichero}{Factor escala}{Caption (leyenda)}{Label (identificador para referenciarla)}
% -------------------------------------------------------------------------------------------------------------
\def\imagen#1#2#3#4{
 \begin{figure}[h]
 \begin{center}
   \scalebox{#2}{\includegraphics{#1}}
 \caption {#3}
 \label{#4}
 \end{center}
 \end{figure}
}


\begin{document}

% En las p�ginas de portada e �ndices, no hay encabezado ni pie de p�gina
\pagestyle{empty} 

% Se incluye la portada
\thispagestyle{empty}
\begin{center}
  {\LARGE UNIVERSIDAD DE CASTILLA-LA MANCHA} \\
  \bigskip
  {\Large ESCUELA SUPERIOR DE INFORM�TICA} \\
  \vspace{28mm}
  \includegraphics[scale=0.45, keepaspectratio]{./imagenes/esi_bw.png} \\
  \vspace{30mm}
  {\LARGE \textbf{Modelos Avanzados de Base de Datos}} \\
  \vspace{14mm}
  {\Large \textsf{\textsc{Listado de preguntas del trabajo Distribuci�n 1:}}} \\
  \vspace{2mm}
  {\Large \textsf{\textsc{Bases de Datos Distribuidas}}}\\
  \vspace{15mm}
  {\large Juan Andrada Romero} \\
  {\large Juan Jos� Antequera Flores} \\
  {\large Jose Domingo L�pez L�pez} \\
  \vspace{9mm}
  {\large \today}
\end{center}
\clearpage

% Texto del reverso de la portada
%%\mbox{}
%%\vspace{18cm}
%%\begin{small}
% Se ajusta la separaci�n entre p�rrafos
%%\parskip=10pt 

%%\copyright~ 2008/2009 Juan Andrada Romero. Universidad de Castilla La Mancha, Escuela Superior de Inform�tica de Ciudad Real.

%%Se permite la modificaci�n, copia y distribuci�n de este documento, seg�n la licencia de documentaci�n GNU (\url{http://www.gnu.org}).

%%Este documento fue compuesto con \LaTeX{}. Im�genes generadas con Power Point y Gimp.
%%\end{small}

%%\newpage

% En las p�ginas de �ndices y prefacio, se utiliza numeraci�n romana
\pagenumbering{Roman}

% Se crea el �ndice
\tableofcontents
% Se pasa p�gina y se a�ade el �ndice de figuras al �ndice principal
%\clearpage\phantomsection
%\addcontentsline{toc}{section}{\listfigurename}
% Se crea el �ndice de figuras
%\listoffigures
% Si se quiere crear un �ndice de tablas se pondr�a: \listoftables

\newpage

% Se ajusta la separaci�n entre p�rrafos
\parskip=10pt

% Se a�ade el prefacio al �ndice
%\clearpage\phantomsection
%\addcontentsline{toc}{section}{Prefacio}

% Comienza el contenido del documento. Se utilizan n�meros ar�bigos y el encabezado y pie de p�gina personalizado
\pagenumbering{arabic}
\pagestyle{fancy}

\section{Bases de datos Web}

\subsection{Introducci�n}

La World Wide Web (Web, de manera abreviada) aparece a finales de los a�os 90, proporcionando una t�cnica simple para explorar grandes vol�menes de p�ginas de informaci�n que residen en Internet. La informaci�n que puede contener dichas p�ginas puede ser texto, gr�ficos, im�genes, sonidos y video, adem�s de los hiperv�nculos que una p�gina web posee hacia otra p�gina. \\
\indent Por todo ello, la Web puede definirse como \textit{un sistema basado en tecnolog�a hipermedia que proporciona un m�todo de explorar informaci�n en Internet de forma no secuencial, utilizando hiperv�nculos}. \cite{libro-dedales}

La Web est� compuesta por una red muy numerosa de computadores, donde unos act�an como \textit{servidores}, para proporcionar la informaci�n, y otros como \textit{clientes} o \textit{navegadores Web}, que solicitan a los servidores la informaci�n, utilizando un protocolo de comunicaci�n. El protocolo que se utiliza para transferir las p�ginas Web a trav�s de Internet es \textbf{HTTP}, que es un protocolo gen�rico, orientado a objetos y sin memoria del estado, es decir, que el servidor no almacena ninguna informaci�n entre una solicitud y la siguiente. Sin embargo, esta carencia de memoria de HTTP hace que sea dif�cil soportar el concepto de \textit{sesi�n}, necesario para poder realizar las operaciones necesarias sobre el SGBD (Sistema Gestor de Bases de Datos). \\
\indent Por otra parte, para componer y poder publicar un documento en la Web, se utiliza el lenguaje HTML, el cu�l define la informaci�n que se transmite entre nodos de la red. Dicho lenguaje ha ido evolucionando a lo largo de su historia, utiliz�ndose actualmente la versi�n HTML 4.01, que tiene soporte para hojas de estilo, scripts y objetos incrustados.

Con todo ello, se puede resumir el funcionamiento de un navegador Web en los siguientes puntos: 
\begin{enumerate}
	\item El usuario, haciendo uso del navegador web (sistema cliente), solicita visualizar una p�gina web.
	\item El navegador solicita esta p�gina (que est� codificada en lenguaje HTML, el lenguaje que puede interpretar) al servidor web (sistema servidor), utilizando el protocolo HTTP.
	\item El servidor web env�a la informaci�n que le han pedido, en HTML, al navegador del usuario.
	\item El navegador muestra la p�gina web en la pantalla del usuario.
\end{enumerate}

Como se puede comprobar, el intercambio de informaci�n entre el cliente y el servidor es posible gracias a que, entre otras cosas, las p�ginas Web utilizan el lenguaje HTML. Sin emabargo, tradicionalmente, se han utilizado documentos HTML est�ticos, para los cuales se creaban las posibles respuestas ante las peticiones del cliente. Pero, en la actualidad, surge la posibilidad de utilizar aplicaciones que permitan acceder a informaci�n de forma din�mica, con contenidos y formatos muy diversos. Esto ha hecho que sea necesaria la integraci�n de los SGBDs en los servidores Web, para poder recuperar la informaci�n de la base de datos seg�n la petici�n que envie el cliente.

As�, en la secci�n \ref{estatica-dinamica} se van a comentar las diferencias existentes entre una p�gina Web est�tica y una p�gina Web din�mica, adem�s de comentar brevemente los servicios Web. En la secci�n \ref{requisitos} se van a comentar las caracter�sticas o requisitos que debe tener una base de datos para su uso en la Web, as� como sus ventajas y desventajas. Para terminar ......... (TXEDO).

\subsection{P�ginas Web est�ticas y din�micas} \label{estatica-dinamica}

Una p�gina web est�tica se puede definir como un documento HTML que contiene la informaci�n sobre la p�gina y donde dicha informaci�n no cambia, a no ser que se cambie el propio documento. Por esta raz�n, no es necesario acceder a una base de datos para recuperar la informaci�n, ya que se �sta se encuentra contenida dentro del propio documento que se va a enviar al cliente. 

Sin embargo, la informaci�n de una p�gina Web din�mica se genera cada vez que se accede a dicha p�gina, realizando, por ejemplo, consultas a una base de datos. De este modo, las p�ginas Web din�micas tienen caracter�sticas que no tienen las p�ginas Web est�ticas:

\begin{milista}
	\item Los datos que se muestran en la p�gina Web son personalizados para cada usuario.
	\item Se puede responder a los datos de entrada que el usuario introduce a trav�s del navegador. 
\end{milista}

En la actualidad, pr�cticamente todas las p�ginas que existen son din�micas, pero existe otro paradigma que est� cobrando una gran imporantica a la hora de construir y utilizar aplicaciones Web: \textbf{Servicios Web}. Un servicio Web puede definirse como \textit{un conjunto de aplicaciones o de tecnolog�as con capacidad para interoperar en la Web, intercambiando datos entre s� con el objetivo de ofrecer unos servicios}.  \cite{servicios-web} En la Figura \ref{fig:servicio} se puede observar un ejemplo de servicios Web.

\imagen{./imagenes/servicioWeb}{0.6}{Interacci�n entre Servicios Web}{fig:servicio}

Las tecnolog�as est�ndares que utiliza un servicio Web para esta interoperabilidad son: 

\begin{milista}
	\item Lenguaje XML (eXtensible Markup Language).
	\item Protocolo SOAP (Simple Object Access Protocol), para la comunicaci�n en Internet.
	\item Lenguaje WSDL (Web Services Description Language), para describir el servicio Web.
	\item Protocolo UDDI (Universal Discovery, Description and Integration), que se utiliza para registrar el servicio y que pueda ser utilizado.
\end{milista}

En el caso de los servicios Web, como al fin y al cabo son aplicaciones Web, es tambi�n necesario el uso de base de datos, para poder acceder a la informaci�n que necesita la aplicaci�n pasa su funcionamiento. 

\subsection{Integraci�n Web-SGBD}

En la Web, no se puede hablar de un tipo de SGBD espec�fico, como puede ocurrir en otros tipos de bases de datos, como pueden ser las orientadas a objetos. Si bien es cierto que existen fabricantes que tratan de proporcionar soluciones propietarias de conectividad, la mayor�a de empresas prefieren utilizar un SGBD m�s general, para no atarse a una tecnolog�a concreta. As�, las bases de datos Web no son m�s que SGBD utilizados y orientados a la Web, de tal modo que se utiliza el SGBD que m�s se adapte a la naturaleza de la informaci�n que debe manejar la aplicaci�n Web correspondiente.  

Los requisitos que debe cumplir un SGBD para poder integrarse en una aplicaci�n Web, son los siguientes: 

\begin{milista}
	\item Incorporar alg�n tipo de \textit{middleware} (capa de software que se sit�a sobre el SGBD) para a�adir conectividad a la base de datos a trav�s de Internet. Dicha conectividad debe ser independiente de los datos y del fabricante.
	\item Soporte para transacciones que abarquen m�ltiples solicitudes HTTP (es decir, acceso concurrente a los datos).
  \item Soporte para mecanismos de autenticaci�n basados en la sesi�n y en la aplicaci�n.
	\item Arquitectura abierta que permita la interoperabilidad con diversos sistemas y tecnolog�as (servicios Web, RMI, CORBA, XML, etc.)
\end{milista}

No obstante, existe un tipo especial de SGBD que se utiliza para datos semi-estructurados, que son datos que pueden ser irregulares o incompletos, cuya estructura puede cambiar de forma r�pida o impredecible y son d�bilmente tipados. \cite{semi-estructurado}. \\
\indent Por este motivo, un SGBD como puede ser un SGBD relacional u orientado a objetos pueden tener problemas a la hora de manejar datos de este tipo, ya que dichos SGBD almacenan los datos siguiendo un esquema fijo, como puede ser una tabla, en el modelo relacional. 

Por tanto, para tratar con este tipo de datos, que abunda en la Web, se utilizan los siguiente modelos y SGBD espec�ficos:

\subsubsection{Modelo de Intercambio de Objetos (OEM - Object Exchange Model}

OEM (\cite{oem}) es un modelo de objetos anidados que fue dise�ado para su uso en diferentes proyectos del \textit{Grupo de Bases de Datos de la Universidad de Standford}, para soportar la integraci�n de datos de diferentes or�genes de datos. Los datos OEM son carentes de esquema y pueden considerarse como un grafo dirigido etiquetado (ver Figura \ref{fig:oem}), donde los nodos son objetos. A su vez, cada objeto est� compuesto por un identificador, una etiqueta descriptiva, un tipo de datos y un valor. En cuanto a los tipos de datos, �stos pueden ser simples (\textit{string}, \textit{int}, etc.) y no tienen conexiones con otros nodos del grafo, o puden ser otros objetos, que se conectan con otros nodos del grafo.

\imagen{./imagenes/OEM}{0.5}{Ejemplo de datos en el modelo OEM}{fig:oem}

\subsubsection{XML - eXtensible Markup Language}

XML (\cite{xml}) es un metalenguaje definido por el World Wide Web Consortium (W3C) que permite la definici�n de etiquetas personalizadas para proporcionar funcionalidad no disponible en HTML. En XML existen dos tipos de modelos: el modelo centrado en los datos y el modelo centrado en los documentos.

Si se utiliza como modelo de documentos, XML se utiliza como formato de almacenamiento e intercambio para datos que estructurados. En este caso, los datos podr�an almacenarse en un SGBD relacional, transformando el documento XML en el esquema relacional. Esto puede hacer con una extensi�n del SQL, llamada SQL/XML.

En el caso de que XML se use para codificar datos semi-estructurados, se usa una base de datos XML nativa (\textit{NXD, Native XML Database}), las cu�les definen un modelo de datos para un documento XML, almacenando y extrayendo documentos de acuerdo con dicho modelo. Ejemplos de estos SGBD son Lore, Infonyte DB o Tamino XML Server. Adem�s, se pueden distinguir dos tipos de NXD, seg�n como almacenen la informaci�n: 
\begin{milista}
	\item \textbf{Basados en texto:} almacenan el c�digo XML como texto, como un archivo de un sistema de archivos o como un dato de tipo CLOB en un SGBD relacional.
	\item \textbf{Basados en modelo:} almacenan el c�digo XML en alguna representaci�n interna en forma de �rbol.
\end{milista}

Por otra parte, estas bases de datos XML utilizan un lenguaje de consulta de datos llamado \textit{XQuery}. Un ejemplo de consulta para obtener una lista de identificadores de empleados de un documento XML, es el siguiente: 
\begin{verbatim}
FOR $S IN doc(�staff_list.xml�)
	WHERE $S/SALARY >15000 AND $S/@branchNo=�B005�
RETURN $S/STAFFNO
\end{verbatim}





\subsection{T�cnicas para la integraci�n Web-SGBD}

En las siguientes subsecciones se comentar�n algunas de las t�cnicas y tecnolog�as m�s utilizadas para la integraci�n de bases de datos en entornos web.

\subsubsection{CGI - Common Gateway Interface}

La interfaz de pasarela com�n es un protocolo est�ndar que define c�mo los \textit{scripts CGI} se comunican con el servidor Web. Un script CGI est� dise�ado para aceptar y devolvder datos que sean conformes con la especificaci�n de CGI, pudiendo as� reutilizar los scripts independientemente del servidor Web que se est� utilizando. 

A continuaci�n se describe la forma de actuaci�n de un CGI de forma esquem�tica:

\begin{enumerate}
	\item En primera instancia, el servidor recibe una petici�n (el cliente ha activado un URL que contiene el CGI), y comprueba si se trata de una invocaci�n de un CGI.
	\item El servidor comprueba los archivos de configuraci�n y de acceso para asegurarse de que el solicitante tiene acceso al script CGI. 
	\item El servidor prepara el entorno para ejecutar el script.
	\item El servidor ejecuta el script, leyendo las variables de entorno y su entrada est�ndar.
	\item El script realiza su funci�n. Como consecuencia de su actividad se va generando un objeto MIME que la aplicaci�n escribe en su salida est�ndar.
	\item Cuando el script finaliza, el servidor env�a la informaci�n producida, junto con informaci�n propia, al cliente, que se encontraba en estado de espera.
\end{enumerate}

Actualmente, �sta es la soluci�n que m�s se est� utilizando para la creaci�n de interfaces Web-SGBD. Entre las ventajas de la programaci�n CGI, se pueden destacar su sencillez, ya que es muy f�cil de entender; es un lenguaje de programaci�n independiente, soportando lenguajes como C, C++, Java, Perl, Python, etc. y es un est�ndar para usarse en todos los servidores Web.

Como desventaja se puede destacar la baja eficiencia que tiene el uso de CGI, debido a que el servidor Web tiene que cargar el script CGI y establecer la conexi�n con la base de datos cada vez que se recibe una petici�n.

\subsubsection{Lenguajes de script}

Un \textit{script} es un c�digo escrito en un lenguaje de porgramaci�n que se incrusta en el c�digo HTML de una p�gina Web y que se puede ejecutar en el lado del servidor o del cliente, a�adiendo nuevas funcionalidades a la p�gina Web. De esta forma, se pueden automatizar funciones para acceder a los datos de la base de datos, creando nuevo c�digo HTML y mostrando as� l�s p�ginas Web personalizadas con los datos que pide el cliente. 

Algunos de los lenguajes de script m�s utilizados actualmente se comentan a continuaci�n.

\paragraph{JavaScript}

JavaScript es un lenguaje de scripts basado en Java, orientado a objetos, interpertado, d�bilmente tipado e incrustado en c�digo HTML. Permite incluir en las p�ginas HTML funciones que respondan a los sucesos del usuario, como clics de rat�n, eventos de teclado, etc., o acceso a la base de datos si se ejecuta en el lado del servidor. 

\paragraph{VBScript}

VBScript es otro lenguaje de scripts procedimental (basado en subrutinas), muy similar a JScript en cuanto a funcionalidad se refiere, pero con la sintaxis de Visual Basic. 

\paragraph{Perl y PHP}

Perl (Practical Extraction and Report Language) es un lenguaje de programaci�n de alto nivel interpretado que combina caracter�sticas de C y de diferentes �rdenes de sistemas UNIX. Hoy en d�a es uno de los lenguajes de script m�s utilizados para su ejecuci�n en el lado del servidor.

Por otra parte, PHP (Hypertext Preprocessor) es otro lenguaje de script de c�digo abierto que se incrusta en HTML y que es soportado por multitud de servidores Web. Uno de los mayores usos que tiene es el acceso a bases de datos en el lado del servidor, utilizando como SGBD \textit{MySQL} y como servidor, \textit{HTTP Apache}. 

\subsubsection{Java}

Java es un lenguaje de programaci�n, orientado a objetos e interpretado desarrollado por Sun Microsystems, que se ha convertido en un est�ndar \textit{de facto} para la programaci�n basada en la Web, permitiendo crear tanto aplicaciones Web (\textbf{applets)}, como aplicaciones de servidor (\textbf{servlets)}. 

Adem�s, Java es un lenguaje independiente de la plataforma, ya que es interpretado en su m�quina virtual (JVM - Java Virtual Machine). Otras caracter�sticas de este lenguaje es que posee un recolector de basura, mecanismos de seguridad con comprobaci�n de tipos y la ausencia de punteros. 

En 1999, Sun crea \textbf{J2EE}, una arquitectura pensada para aplicaciones empresariales robustas, escalables, multiusuario y seguras. Esta plataforma incluye todo lo necesario para el desarrollo de aplicaciones Web, aplicaciones de servidor, presentaci�n en el cliente y las comunicaciones necesarias para comunicar el cliente, el servidor y la bases de datos (ver Figura \ref{fig:j2ee}). 

\imagen{./imagenes/j2ee}{0.5}{Arquitectura J2EE simplificada}{fig:j2ee}

\paragraph{JDBC}

JDBC es la t�cnica mas utilizada y madura para el acceso a un SGBD relacional desde Java. JDBC define una API para acceder a diferentes SGBD relacionales, soportando toda la funcionalidad SQL b�sica. Adem�s, por encima del JDBC se construyen interfaces de aplicaci�n a mayor nivel como,  por ejemplo, SQL embebido en Java o una correspondencia directa de tablas de una base de datos relacional a clases Java (Hibernate). 

Para el acceso a la base de datos por parte de las aplicaciones, hay diferentes opciones:

\begin{milista}
	\item Usar el puente JDBC-ODBC, donde el controlador ODBC act�a como nivel mediador entre el SGBD y el controlador JDBC. Es necesario cuando no se puede acceder directamente a las bibliotecas del fabricante del SGBD.
	\item Usar el controlador JDBC parcial, que convierte llamadas JDBC directamente a llamadas del API del SGBD.
	\item Usar JDBC puro para middleware de bases de datos, que traduce las llamadas del controlador JDBC al protocolo del middleware, el cual es traducido de nuevo al protocolo del SGBD. De este modo, existe m�s variedad a la hora de usar diferentes SGBD.
	\item Usar JDBC puro con conexi�n directa al SGBD. Esto conlleva utilizar un controlador diferente, seg�n el SGBD que se utilice. 
\end{milista}

Sea cual sea la opci�n de comunicaci�n elegida, el proceso de interacci�n con la base de datos es el siguiente: 

\begin{enumerate}
	\item Registrarse en el controlador JDBC, utilizando el m�todo \texttt{Class.forName()}.
	\item Abrir la conexi�n con el SGBD, utilizando el m�todo \texttt{DriverManager.getConnection()}, indicando la URL del SGBD, que tiene el siguiente formato: \texttt{protocolo:subprotocolo:subnombre}. Por ejemplo: \texttt{jdbc:odbc:www.dreamhome.es.co.uk/dhdatabase}.
	\item Ejecutar una sentencia SQL sobre la base de datos, con los m�todo \texttt{createStatement} y texttt{execute()}.
	\item Procesar los resultados obtenidos (si la sentencia era una consulta).
	\item Cerrar la conexi�n, con el m�todo \texttt{close()}.
\end{enumerate}



% A�adimos la bibliografia al �ndice
\clearpage\phantomsection
\addcontentsline{toc}{section}{\bibname}

% Bibliograf�a. En este caso se usa BibTeX
\bibliographystyle{plain}
\pagestyle{plain} 
\bibliography{Bibliografia}

\end{document}
